%==============================================================================
%  MATH 231 -- Linear Algebra I (Queens College, Fall 2026)
%  MIDTERM 1, PRACTICE #2 -- built from Homework 1 and Homework 2
%
%  COMPANION TO PRACTICE #1.  Same scope, same five parts, and the same number
%  of questions in each, laid out page for page the same way.  The values are
%  new throughout, the true-or-false and multiple-choice statements are new,
%  and Part E asks for different proofs.
%
%  SCOPE.  Exactly the material of the two homework assignments:
%     homework/MATH231_homework_1.tex   Parts A-E (Part F, NumPy, left out)
%     homework/MATH231_homework_2.tex   Parts A-F
%  Programming is not assessed on exams (syllabus S7), so there is no NumPy.
%  Every question is a close relative of a homework problem with new numbers,
%  not a verbatim copy -- the syllabus says exam questions are not recycled from
%  homework.  The table on the last page maps each question to its sources.
%
%  FORMAT.  Five parts, by question type:
%     A  true or false        A.1-A.10
%     B  multiple choice      B.1-B.8
%     C  computation          C.1-C.6
%     D  applications         D.1-D.4
%     E  short proofs         E.1-E.5
%  This is longer than one class period, on purpose: it is meant to touch every
%  topic of both homework sets.  Trim Parts C and D first if a timed version is
%  wanted.
%
%  SPACING.  Items on a page are separated by \vfill, so the blank space on each
%  page is shared out evenly below every question; each page ends with an
%  explicit \newpage.  A problem that runs past one page continues under a
%  "Problem X.n, continued" line.  Moving an item means checking the page breaks
%  again.
%
%  CALCULATOR-FREE.  Every answer is exact and clean: angles are 0, 30, 45, 60,
%  90, 120 or 135 degrees, comparisons of square roots can be made by squaring,
%  and no part asks for a decimal.  There is no answer key in this file.
%
%  Compile with tectonic (or pdflatex; standard TeX Live packages only).
%==============================================================================

\documentclass[11pt]{article}

\usepackage[margin=1in]{geometry}
\usepackage{amsmath,amssymb}
\usepackage[dvipsnames]{xcolor}
\usepackage{enumitem}
\usepackage{booktabs}
\usepackage{array}
\usepackage[hidelinks]{hyperref}
\usepackage{titlesec}
\usepackage{needspace}

%------------------------------------------------------------------ style ----
\titleformat{\section}{\large\bfseries\sffamily\color{RoyalBlue}}
                      {Part \thesection.}{0.6em}{}
\titlespacing*{\section}{0pt}{0pt}{7pt}

\renewcommand{\thesection}{\Alph{section}}

% Keeps long inline formulas from pushing a line into the margin.
\setlength{\emergencystretch}{3em}

% The problem heading is a line of its own, so that part (a) -- or the opening
% sentence -- always starts below it.  Same environment as Homework 2.
\makeatletter
\newcounter{prob}[section]
\renewcommand{\theprob}{\thesection.\arabic{prob}}
\newenvironment{prob}[1][]
  {\par\addvspace{12pt plus 4pt minus 4pt}%
   \needspace{4\baselineskip}%
   \refstepcounter{prob}%
   \noindent\textbf{Problem~\theprob}%
   \if\relax\detokenize{#1}\relax\else\ (#1)\fi\textbf{.}%
   \par\nobreak\@afterindentfalse\@afterheading}
  {\par}
\makeatother

\newcommand{\hint}[1]{\par\smallskip\noindent
  {\small\sffamily\textcolor{RoyalBlue}{\textbf{Hint.}}\ #1}\par\smallskip}

% A problem that runs onto a second page continues under this line.
\newcommand{\contd}[1]{\newpage\noindent
  {\small\sffamily\color{gray}Problem #1, continued}\par\medskip}

% True or false: the statement on the left, the two letters to circle on the
% right of its first line.
\newcommand{\tf}[1]{\parbox[t]{\dimexpr\linewidth-5.5em}{#1}%
  \hfill\textsf{\textbf{T}\qquad\textbf{F}}}

% Multiple choice: four short options in one row, medium ones two by two, or
% long ones one per line.
\newcommand{\fourch}[4]{\par\smallskip\noindent
  \begin{tabular}{@{}*{4}{>{\raggedright\arraybackslash}p{0.22\linewidth}}@{}}
  (A)\ #1 & (B)\ #2 & (C)\ #3 & (D)\ #4
  \end{tabular}\par}
\newcommand{\twoch}[4]{\par\smallskip\noindent
  \begin{tabular}{@{}*{2}{>{\raggedright\arraybackslash}p{0.46\linewidth}}@{}}
  (A)\ #1 & (B)\ #2 \\[6pt]
  (C)\ #3 & (D)\ #4
  \end{tabular}\par}
\newcommand{\stackch}[4]{\par\smallskip\noindent
  \begin{tabular}{@{}l@{}}
  (A)\ #1 \\[3pt] (B)\ #2 \\[3pt] (C)\ #3 \\[3pt] (D)\ #4
  \end{tabular}\par}

% shaded call-out box.  The body is collected into a savebox first, because
% \colorbox needs its whole argument at once.
\definecolor{softbg}{HTML}{F2F6FA}
\newsavebox{\qcbox}
\newenvironment{callout}[1]
  {\par\medskip\needspace{5\baselineskip}%
   \begin{lrbox}{\qcbox}%
   \begin{minipage}{\dimexpr\textwidth-2\fboxsep-4pt}%
   \setlength{\parskip}{0.4em}%
   \textbf{\sffamily\color{RoyalBlue}#1}\par}
  {\end{minipage}\end{lrbox}%
   \begin{center}\colorbox{softbg}{\usebox{\qcbox}}\end{center}\medskip}

% list styles used throughout
\newlist{tfitems}{enumerate}{1}
\setlist[tfitems]{leftmargin=3.2em,label=\textbf{\thesection.\arabic*},
                  itemsep=0pt,topsep=4pt}
\newlist{parts}{enumerate}{1}
\setlist[parts]{leftmargin=2.6em,label=\textbf{(\alph*)},itemsep=0pt,
                topsep=6pt}
%--------------------------------------------------------------- notation ----
\newcommand{\R}{\mathbb{R}}
\newcommand{\vc}[1]{\mathbf{#1}}
\newcommand{\uv}[1]{\hat{\vc{#1}}}
\newcommand{\one}{\mathbf{1}}
\newcommand{\norm}[1]{\lVert #1\rVert}
\newcommand{\ip}[2]{\langle #1,#2\rangle}
\newcommand{\proj}[2]{\operatorname{proj}_{#1}(#2)}
\newcommand{\comp}[2]{\operatorname{comp}_{#1}(#2)}
\newcommand{\conj}[1]{\overline{#1}}
\newcommand{\Rez}{\operatorname{Re}}
\newcommand{\Imz}{\operatorname{Im}}
\DeclareMathOperator{\spanop}{span}
\newcommand{\grad}{\nabla}
\newcommand{\pd}[2]{\frac{\partial #1}{\partial #2}}

\title{\vspace{-1.2cm}\textbf{\sffamily MATH 231 -- Linear Algebra I}\\[4pt]
       {\large\sffamily Midterm 1 Practice \#2 \textbullet\ Homework 1 and
        Homework 2 material}\\[2pt]
       {\normalsize\sffamily Kiely Hall 326 \textbullet\
        Instructor: Kevin Bedoya}}
\author{}
\date{}

\begin{document}
\maketitle
\vspace{-1.4cm}

\begin{center}
\large\sffamily Name: \rule{9cm}{0.4pt}
\end{center}

\begin{callout}{About this practice exam}
Every question here is a close relative of a problem from Homework~1 or
Homework~2, with new numbers. The table on the last page lists the homework
problems each question is modeled on, so you know where to look when one gives
you trouble.

This is practice. It is not a statement about the format, the length, or the
questions of the actual midterm.
\end{callout}

\begin{callout}{Instructions}
\begin{itemize}[leftmargin=1.4em,itemsep=5pt,topsep=3pt]
  \item \textbf{Parts A and B.} Circle one answer. No justification is needed;
        the space below each question is for scratch work.
  \item \textbf{Parts C, D and E.} Show your work in the space below each
        question. A correct answer with no derivation earns partial credit at
        best. If you run out of room, continue on the back of the page and say
        so.
  \item Leave answers in exact form: $\sqrt{13}$, $\pi/3$, $\ln2$,
        $-\tfrac45+\tfrac75 i$. No question needs a calculator.
  \item \textbf{Proofs.} Write down what you are assuming, justify each
        equality by naming the fact that licenses it, and end with a sentence
        saying what you have shown. A chain of equalities with no words is not a
        proof.
\end{itemize}
\end{callout}

\begin{center}
\renewcommand{\arraystretch}{1.3}
\begin{tabular}{@{}c l l@{}}
\toprule
\textbf{Part} & \textbf{Type} & \textbf{Questions} \\
\midrule
A & True or false   & A.1--A.10 \\
B & Multiple choice & B.1--B.8 \\
C & Computation     & C.1--C.6 \\
D & Applications    & D.1--D.4 \\
E & Short proofs    & E.1--E.5 \\
\bottomrule
\end{tabular}
\end{center}

\newpage
%==============================================================================
\begin{callout}{Notation used throughout}
Bold letters ($\vc{u}$, $\vc{v}$, $\vc{x}$) are vectors; plain letters
($\alpha$, $k$, $t$) are scalars, that is, real numbers. For
$\vc{u}=[\,u_1,\dots,u_n\,]$ and $\vc{v}=[\,v_1,\dots,v_n\,]$ in $\R^{n}$:
\[
  \ip{\vc{u}}{\vc{v}}=\vc{u}\cdot\vc{v}=\sum_{i=1}^{n}u_iv_i,
  \qquad
  \norm{\vc{v}}=\norm{\vc{v}}_2=\sqrt{\ip{\vc{v}}{\vc{v}}},
  \qquad
  \norm{\vc{v}}_1=\sum_{i=1}^{n}\lvert v_i\rvert,
  \qquad
  \norm{\vc{v}}_\infty=\max_{1\le i\le n}\lvert v_i\rvert ,
\]
and each norm gives a distance $d_p(\vc{x},\vc{y})=\norm{\vc{x}-\vc{y}}_p$.

For nonzero $\vc{u},\vc{v}$, the \textbf{cosine similarity} is
$\cos(\vc{u},\vc{v})=\dfrac{\ip{\vc{u}}{\vc{v}}}{\norm{\vc{u}}\,\norm{\vc{v}}}$,
and the angle $\theta\in[0,\pi]$ between $\vc{u}$ and $\vc{v}$ is the one with
$\cos\theta=\cos(\vc{u},\vc{v})$. The \textbf{unit vector} in the direction of
$\vc{u}$ is $\uv{u}=\vc{u}/\norm{\vc{u}}$.

For $\vc{v}\neq\vc{0}$, the \textbf{component} of $\vc{u}$ along $\vc{v}$ (a
scalar) and the \textbf{projection} of $\vc{u}$ onto $\vc{v}$ (a vector) are
\[
  \comp{\vc{v}}{\vc{u}}=\frac{\ip{\vc{v}}{\vc{u}}}{\norm{\vc{v}}},
  \qquad
  \proj{\vc{v}}{\vc{u}}=\frac{\ip{\vc{v}}{\vc{u}}}{\norm{\vc{v}}^{2}}\,\vc{v},
\]
and $\vc{u}-\proj{\vc{v}}{\vc{u}}$ is the \textbf{orthogonal piece}.

In $\R^{n}$, $\vc{e}_i$ is the standard basis vector with a $1$ in slot $i$ and
$0$ elsewhere, and $\one=[\,1,\,1,\,\dots,\,1\,]$ is the all-ones vector.

A \textbf{subspace of $\R^{2}$} is a subset of $\R^{2}$ that is non-empty, closed
under addition, and closed under scalar multiplication.

For a complex number $z=a+bi$ with $a,b\in\R$: $\Rez(z)=a$, $\Imz(z)=b$,
$\conj{z}=a-bi$, $\lvert z\rvert=\sqrt{a^{2}+b^{2}}$, and Euler's formula is
$e^{i\theta}=\cos\theta+i\sin\theta$.

For a function $f(x,y)$, the \textbf{gradient} is
$\grad f=\bigl[\,\partial f/\partial x,\ \partial f/\partial y\,\bigr]$.
\end{callout}

\newpage
%==============================================================================
\section{True or false}
%==============================================================================
Circle \textbf{T} if the statement is always true, and \textbf{F} otherwise.

\begin{tfitems}
  \item \tf{For all vectors $\vc{u},\vc{v}\in\R^{2}$,
        $\norm{\vc{u}-\vc{v}}=\norm{\vc{v}-\vc{u}}$.}
        \vfill
  \item \tf{For all complex numbers $z$ and $w$,
        $\lvert z+w\rvert=\lvert z\rvert+\lvert w\rvert$.}
        \vfill
  \item \tf{$i^{99}=-i$.}
        \vfill
  \item \tf{If $\vc{u}\in\R^{2}$ is nonzero and $\alpha$ is a nonzero scalar,
        then $\proj{\alpha\vc{u}}{\vc{v}}=\proj{\vc{u}}{\vc{v}}$ for every
        $\vc{v}\in\R^{2}$.}
        \vfill
  \item \tf{For all nonzero $\vc{u},\vc{v}\in\R^{2}$,
        $\cos(-\vc{u},\vc{v})=\cos(\vc{u},\vc{v})$.}
        \vfill
\end{tfitems}

\newpage
\begin{tfitems}[start=6]
  \item \tf{For every $\vc{v}\in\R^{2}$,
        $\norm{\vc{v}}_1\le2\,\norm{\vc{v}}_\infty$.}
        \vfill
  \item \tf{The set $\{(x,y) : x\ge0 \text{ and } y\ge0\}$ is a subspace of
        $\R^{2}$.}
        \vfill
  \item \tf{If $\vc{u},\vc{v}\in\R^{2}$ are linearly independent, then
        $\vc{u}+\vc{v}$ and $\vc{u}-\vc{v}$ are linearly independent.}
        \vfill
  \item \tf{If $\vc{u},\vc{v}\in\R^{2}$ are nonzero and
        $\ip{\vc{u}}{\vc{v}}\neq0$, then $\vc{u}$ and $\vc{v}$ are linearly
        dependent.}
        \vfill
  \item \tf{If $f(x,y)=x\,e^{y}$, then $\grad f(1,0)=[\,1,\,0\,]$.}
        \vfill
\end{tfitems}

\newpage
%==============================================================================
\section{Multiple choice}
%==============================================================================
Circle the one correct answer.

\begin{tfitems}
  \item Which vector has length $15$ and points in the same direction as
        $[\,-3,\,4\,]$?
        \fourch{$[\,-45,\,60\,]$}
               {$[\,9,\,-12\,]$}
               {$[\,-9,\,12\,]$}
               {$[\,-\tfrac35,\,\tfrac45\,]$}
        \vfill
  \item The quotient $\dfrac{4+3i}{2-i}$ equals
        \fourch{$2+i$}
               {$1+2i$}
               {$1-2i$}
               {$5+10i$}
        \vfill
  \item The set of complex numbers $z$ with $\lvert z+2i\rvert\le1$ is
        \stackch{the circle of radius $1$ centred at $-2i$}
              {the closed disc of radius $1$ centred at $2i$}
              {the closed disc of radius $1$ centred at $-2i$}
              {the open disc of radius $1$ centred at $-2i$}
        \vfill
\end{tfitems}

\newpage
\begin{tfitems}[start=4]
  \item The angle between $[\,2,\,0\,]$ and $[\,-1,\,\sqrt3\,]$ is
        \fourch{$\pi/6$}{$\pi/3$}{$\pi/2$}{$2\pi/3$}
        \vfill
  \item For $\vc{u}=[\,2,\,3\,]$ and $\vc{v}=[\,3,\,4\,]$, the component
        $\comp{\vc{v}}{\vc{u}}$ is
        \fourch{$\tfrac{18}{5}$}{$\tfrac{18}{25}$}
               {$[\,\tfrac{54}{25},\,\tfrac{72}{25}\,]$}
               {$\tfrac{18}{\sqrt{13}}$}
        \vfill
  \item Which pair is a basis of $\R^{2}$?
        \twoch{$[\,2,\,-4\,]$ and $[\,-1,\,2\,]$}
              {$[\,3,\,6\,]$ and $[\,1,\,2\,]$}
              {$[\,0,\,0\,]$ and $[\,1,\,3\,]$}
              {$[\,1,\,2\,]$ and $[\,2,\,1\,]$}
        \vfill
\end{tfitems}

\newpage
\begin{tfitems}[start=7]
  \item In $\R^{9}$, the norms $\norm{\one}_1$, $\norm{\one}_2$ and
        $\norm{\one}_\infty$ are, in that order,
        \fourch{$1,\ 3,\ 9$}{$9,\ 3,\ 1$}{$9,\ 81,\ 1$}{$9,\ 9,\ 9$}
        \vfill
  \item For $f(x,y)=(x-3y)^{4}$, the partial derivative
        $\partial f/\partial y$ is
        \fourch{$-12(x-3y)^{3}$}{$4(x-3y)^{3}$}{$-3(x-3y)^{3}$}
               {$12(x-3y)^{3}$}
        \vfill
\end{tfitems}

\newpage
%==============================================================================
\section{Computation}
%==============================================================================
Show your work. Leave every answer in exact form.

%------------------------------------------------------------------------------
\begin{prob}[Vectors in the plane]
Let $\vc{u}=[\,-1,\,3\,]$ and $\vc{v}=[\,4,\,1\,]$ be vectors in $\R^{2}$, and
let $\alpha=-3$ be a scalar.
\begin{parts}
  \item Compute $\vc{u}-\vc{v}$, \ $3\vc{u}+2\vc{v}$, \ and $\alpha\vc{u}$,
        showing the componentwise step.
        \vfill
  \item Compute $\norm{\vc{u}}$, $\norm{\vc{v}}$, $\norm{\vc{u}+\vc{v}}$ and
        $\norm{\alpha\vc{u}}$.
        \vfill
  \item For this pair, does the triangle inequality
        $\norm{\vc{u}+\vc{v}}\le\norm{\vc{u}}+\norm{\vc{v}}$ hold strictly
        ($<$) or with equality ($=$)?
        \vfill
\end{parts}
\contd{C.1}
\begin{parts}[resume]
  \item Find every scalar $\beta$ for which
        $\norm{\vc{u}+\beta\vc{u}}=\norm{\vc{u}}+\norm{\beta\vc{u}}$. You may
        use $\norm{\gamma\vc{x}}=\lvert\gamma\rvert\,\norm{\vc{x}}$ for every
        scalar $\gamma$ and every vector $\vc{x}$.
        \vfill
  \item The arrow from $(2,-1)$ to $(1,2)$ and the arrow from $(0,4)$ to $(3,3)$
        both have length $\sqrt{10}$. Write each arrow as a vector, using
        \emph{head minus tail}, and decide which of them, if either, is the
        vector $\vc{u}$. Justify your answer with the definition of equality of
        vectors.
        \vfill
\end{parts}
\end{prob}

\newpage
%------------------------------------------------------------------------------
\begin{prob}[Complex arithmetic]
Let $z=4-i$ and $w=2+3i$. Give every answer that is not a modulus in the form
$a+bi$ with $a$ and $b$ real.
\begin{parts}
  \item Compute $z+w$, \ $z-w$, \ and $zw$.
        \vfill
  \item Compute $\conj{w}$, \ $w\conj{w}$, \ and $\dfrac{z}{w}$.
        \vfill
  \item Compute $\lvert z\rvert$, $\lvert w\rvert$ and $\lvert zw\rvert$, and
        compare $\lvert zw\rvert$ with $\lvert z\rvert\,\lvert w\rvert$.
        \vfill
\end{parts}
\contd{C.2}
\begin{parts}[resume]
  \item Compute $(1-i)^{2}$ and $(1-i)^{6}$.
        \vfill
  \item Use Euler's formula to write $e^{3\pi i/4}$ and $e^{-i\pi/3}$ in the
        form $a+bi$.
        \vfill
  \item The point $(-3,4)$ is rotated $90^{\circ}$ \emph{clockwise} about the
        origin. Find the coordinates of the rotated point by a single
        multiplication of complex numbers.
        \vfill
\end{parts}
\end{prob}

\newpage
%------------------------------------------------------------------------------
\begin{prob}[Dot products, angles, and orthogonality]
For each pair below, compute $\ip{\vc{u}}{\vc{v}}$, $\norm{\vc{u}}$,
$\norm{\vc{v}}$ and $\cos(\vc{u},\vc{v})$. Then give the angle $\theta$ between
the two vectors, and state which of \emph{parallel}, \emph{antiparallel},
\emph{orthogonal} or \emph{none of the three} describes the pair.
\begin{center}
\renewcommand{\arraystretch}{1.25}
\begin{tabular}{@{}c c c @{\qquad\qquad} c c c@{}}
\toprule
 & $\vc{u}$ & $\vc{v}$ & & $\vc{u}$ & $\vc{v}$ \\
\midrule
(a) & $[\,3,\,-2\,]$ & $[\,4,\,6\,]$  & (c) & $[\,2,\,1\,]$       & $[\,-3,\,1\,]$ \\
(b) & $[\,-1,\,2\,]$ & $[\,-3,\,6\,]$ & (d) & $[\,\sqrt3,\,1\,]$  & $[\,0,\,2\,]$  \\
\bottomrule
\end{tabular}
\end{center}
\vfill
\contd{C.3}
\begin{parts}[start=5]
  \item Find every real number $k$ for which $[\,k,\,2\,]$ is orthogonal to
        $[\,k+1,\,-1\,]$.
        \vfill
  \item Find every real number $k$ for which the angle between
        $[\,1,\,k\,]$ and $[\,0,\,1\,]$ is $\pi/6$. You will want to square
        both sides; check every candidate in the original equation, and discard
        any that fail it.
        \vfill
  \item In $\R^{4}$, let $\vc{u}=[\,1,\,-1,\,2,\,0\,]$ and
        $\vc{w}=[\,1,\,0,\,1,\,1\,]$. Compute $\ip{\vc{u}}{\vc{w}}$,
        $\norm{\vc{u}}$, $\norm{\vc{w}}$, and the angle between $\vc{u}$ and
        $\vc{w}$.
        \vfill
\end{parts}
\end{prob}

\newpage
%------------------------------------------------------------------------------
\begin{prob}[Bases, independence, and coordinates]
Two vectors $\vc{u},\vc{v}\in\R^{2}$ are \textbf{linearly independent} when
$\alpha\vc{u}+\beta\vc{v}=\vc{0}$ forces $\alpha=\beta=0$, and
\textbf{linearly dependent} otherwise. For $\vc{u}=[\,u_1,\,u_2\,]$ and
$\vc{v}=[\,v_1,\,v_2\,]$, put $D=u_1v_2-u_2v_1$. The pair is linearly
independent, and so a basis of $\R^{2}$, exactly when $D\neq0$ (Vector Unit,
Part~4, \S18.5).
\begin{parts}
  \item Let $\vc{u}=[\,2,\,1\,]$ and $\vc{v}=[\,-1,\,3\,]$. Compute $D$. Then
        find the coordinates $\alpha,\beta$ of $\vc{w}=[\,3,\,5\,]$ and of
        $\vc{w}=[\,4,\,-5\,]$ against the basis $\{\vc{u},\vc{v}\}$ of
        $\R^{2}$, by setting up the linear system
        $\alpha\vc{u}+\beta\vc{v}=\vc{w}$ and eliminating by hand. Check each
        answer by substituting back.
        \vfill
  \item Show that $[\,-6,\,9\,]$ and $[\,4,\,-6\,]$ are linearly dependent by
        giving scalars $\alpha,\beta$, not both zero, with
        $\alpha[\,-6,\,9\,]+\beta[\,4,\,-6\,]=\vc{0}$. Describe their span as a
        set $\{(x,y) : \dots\}$, give a basis of it, and state its dimension.
        \vfill
\end{parts}
\contd{C.4}
\begin{parts}[resume]
  \item Find every real number $k$ for which $[\,2,\,k\,]$ and $[\,k,\,8\,]$
        are linearly dependent. For each such $k$, give the equation of the line
        through the origin that the two vectors span.
        \vfill
  \item Let $S=\bigl\{[\,1,\,2\,],\ [\,3,\,-1\,],\ [\,4,\,1\,]\bigr\}$. Write
        $[\,5,\,3\,]$ as a linear combination of the three vectors of $S$ in two
        different ways. A basis of $\R^{2}$ must span $\R^{2}$ and be linearly
        independent; $S$ spans $\R^{2}$. Use your two combinations to show which
        of the two conditions $S$ fails.
        \vfill
\end{parts}
\end{prob}

\newpage
%------------------------------------------------------------------------------
\begin{prob}[Orthogonal bases]
If $\vc{b}_1,\dots,\vc{b}_k$ are nonzero, mutually orthogonal, and form a basis,
then the coefficient of $\vc{b}_j$ in a vector $\vc{w}$ is
$\ip{\vc{b}_j}{\vc{w}}/\norm{\vc{b}_j}^{2}$ (Vector Unit, Part~5, \S21.6).
\begin{parts}
  \item Let $\vc{u}=[\,3,\,-1\,]$ and $\vc{v}=[\,1,\,3\,]$. Check that they
        are nonzero and orthogonal. Find the coordinates of $[\,-1,\,7\,]$
        against the basis $\{\vc{u},\vc{v}\}$ of $\R^{2}$ by projection, and
        check them by substituting back.
        \vfill
  \item In $\R^{3}$, let $\vc{u}=[\,1,\,1,\,0\,]$, $\vc{v}=[\,1,\,-1,\,2\,]$
        and $\vc{w}=[\,-1,\,1,\,1\,]$. Check that they are mutually orthogonal.
        Find, by projection, the coefficients $\alpha,\beta,\gamma$ with
        $\vc{x}=\alpha\vc{u}+\beta\vc{v}+\gamma\vc{w}$ for
        $\vc{x}=[\,4,\,0,\,1\,]$, and substitute back to confirm. Then verify
        that
        $\norm{\vc{x}}^{2}
         =\alpha^{2}\norm{\vc{u}}^{2}+\beta^{2}\norm{\vc{v}}^{2}
          +\gamma^{2}\norm{\vc{w}}^{2}$.
        \vfill
\end{parts}
\contd{C.5}
\begin{parts}[resume]
  \item Let $\vc{u}=[\,1,\,1\,]$ and $\vc{v}=[\,3,\,1\,]$. Show that they are
        not orthogonal. Compute $\vc{v}'=\vc{v}-\proj{\vc{u}}{\vc{v}}$ and
        check that $\vc{v}'$ is orthogonal to $\vc{u}$. Find the coordinates of
        $[\,5,\,1\,]$ against the orthogonal basis $\{\vc{u},\vc{v}'\}$ by
        projection. Then rewrite the result as a combination of the original
        $\vc{u}$ and $\vc{v}$, and check it.
        \vfill
\end{parts}
\end{prob}

\newpage
%------------------------------------------------------------------------------
\begin{prob}[Derivatives, partial derivatives, and the gradient]
\begin{parts}
  \item Differentiate each function, and name the rule or rules you use.
        \[
          \text{(i) } x^{3}\ln x
          \qquad
          \text{(ii) } e^{-x}\cos x
          \qquad
          \text{(iii) } \ln\bigl(1+e^{2x}\bigr)
          \qquad
          \text{(iv) } \frac{x}{1+x^{2}}
        \]
        \vfill
  \item Solve $e^{2x}+e^{x}-6=0$.
        \vfill
\end{parts}
\contd{C.6}
\begin{parts}[resume]
  \item Solve $\ln(x+1)+\ln(x-2)=\ln 4$. Reject any candidate that does not
        solve the original equation, and say why it fails.
        \vfill
  \item Compute $\partial f/\partial x$ and $\partial f/\partial y$ for
        \[
          f(x,y)=2x^{4}y-3xy^{3}+5
          \qquad\text{and}\qquad
          f(x,y)=x\,e^{-xy}.
        \]
        \vfill
  \item Let $h(x,y)=x^{2}y-2xy+y^{2}$. Compute $\grad h$, the vector
        $\grad h(0,2)$ and its norm, and find every point $(x,y)$ at which
        $\grad h=\vc{0}$.
        \vfill
\end{parts}
\end{prob}

\newpage
%==============================================================================
\section{Applications}
%==============================================================================
Show your work, and give units wherever the problem has them.

%------------------------------------------------------------------------------
\begin{prob}[Work, and a box on a ramp]
A force is a vector $\vc{F}\in\R^{2}$: its direction is the way it pushes, and
its norm $\norm{\vc{F}}$ is how hard. When a constant force $\vc{F}$ acts on an
object that moves in a straight line with displacement $\vc{s}\in\R^{2}$, the
\textbf{work} done by the force is the scalar $W=\ip{\vc{F}}{\vc{s}}$, in joules
when $\vc{F}$ is in newtons and $\vc{s}$ in metres. So $\vc{F}$, $\vc{s}$,
$\vc{d}$, $\vc{G}$ and $\vc{P}$ below are vectors, while $W$, every norm, and
every component are scalars.

A box sits on a ramp that rises in the direction $\vc{d}=[\,3,\,1\,]$, and
gravity pulls on it with the constant force $\vc{G}=[\,0,\,-20\,]$.
\begin{parts}
  \item Compute $\comp{\vc{d}}{\vc{G}}$, $\proj{\vc{d}}{\vc{G}}$, and the
        orthogonal piece $\vc{G}-\proj{\vc{d}}{\vc{G}}$. Check that the
        orthogonal piece is orthogonal to $\vc{d}$.
        \vfill
  \item Compute the norms of the two pieces. One pulls the box down the ramp
        and the other presses it into the ramp's surface: which is which? Check
        that $\norm{\vc{G}}^{2}$ is the sum of the squares of the two norms.
        \vfill
\end{parts}
\contd{D.1}
\begin{parts}[resume]
  \item The box is pushed up the ramp from $(0,0)$ to $(6,2)$. Compute the
        work done by gravity, and say what its sign means.
        \vfill
  \item Over the same displacement, a worker pushes horizontally with the
        force $\vc{P}=[\,25,\,0\,]$. Compute the work the worker does in two
        ways: as $\ip{\vc{P}}{\vc{s}}$, and as
        $\comp{\vc{s}}{\vc{P}}\cdot\norm{\vc{s}}$. Then give a nonzero force
        that does no work at all over this displacement, and justify it with
        $\ip{\vc{F}}{\vc{s}}=\norm{\vc{F}}\,\norm{\vc{s}}\cos\theta$.
        \vfill
\end{parts}
\end{prob}

\newpage
%------------------------------------------------------------------------------
\begin{prob}[Which customer is nearest? It depends on the metric]
A depot sits at $P=(1,-2)$, and three customers at $A=(1,5)$, $B=(-4,-7)$ and
$C=(3,-8)$.
\begin{parts}
  \item Make a table of the distances $d_1$, $d_2$ and $d_\infty$ from $P$ to
        each customer.
        \vfill
  \item Three vehicles leave the depot: a \emph{taxi} that must follow a
        rectangular street grid, a \emph{drone} that flies in a straight line,
        and an \emph{overhead crane} whose east--west and north--south motors run
        at the same time and at the same speed. Which of $d_1$, $d_2$, $d_\infty$
        measures the trip for each vehicle, and which customer is nearest for
        each?
        \vfill
  \item On one set of axes, sketch and label the three sets
        $\{\vc{v}\in\R^{2} : \norm{\vc{v}}_1=2\}$,
        $\{\vc{v}\in\R^{2} : \norm{\vc{v}}_2=2\}$ and
        $\{\vc{v}\in\R^{2} : \norm{\vc{v}}_\infty=2\}$. For each one, find the
        point where it crosses the ray $y=x$, $x>0$.
        \vfill
\end{parts}
\end{prob}

\newpage
%------------------------------------------------------------------------------
\begin{prob}[Comparing documents]
One way to compare documents is to count words. Fix the vocabulary
\emph{goal}, \emph{team}, \emph{coach}, \emph{vote}, \emph{budget}, and record
each document as the vector in $\R^{5}$ of its word counts, in that order:
\[
  \vc{a}=[\,1,\,2,\,2,\,0,\,0\,],\qquad
  \vc{b}=[\,3,\,6,\,6,\,0,\,0\,],\qquad
  \vc{c}=[\,0,\,0,\,0,\,4,\,3\,],\qquad
  \vc{d}=[\,2,\,1,\,0,\,2,\,0\,].
\]
\begin{parts}
  \item Compute $\cos(\vc{a},\vc{b})$, $\cos(\vc{a},\vc{c})$ and
        $\cos(\vc{a},\vc{d})$.
        \vfill
  \item Compute the Euclidean distances $d_2(\vc{a},\vc{b})$,
        $d_2(\vc{a},\vc{c})$ and $d_2(\vc{a},\vc{d})$.
        \vfill
  \item Rank $\vc{b}$, $\vc{c}$ and $\vc{d}$ from closest to farthest from
        $\vc{a}$, once by cosine similarity and once by Euclidean distance.
        Document $\vc{b}$ is document $\vc{a}$ pasted in three times, and
        $\vc{c}$ shares no words with $\vc{a}$. Which of the two measures gives
        the sensible ranking for comparing what documents are about, and why?
        \vfill
\end{parts}
\end{prob}

\newpage
%------------------------------------------------------------------------------
\begin{prob}[How kinetic energy responds to mass and speed]
An object of mass $m$ kilograms moving at speed $v$ metres per second has
kinetic energy $K(m,v)=\tfrac12 mv^{2}$ joules.
\begin{parts}
  \item Compute $\partial K/\partial m$ and $\partial K/\partial v$.
        \vfill
  \item Evaluate both partial derivatives for a car with $m=1000$ and $v=20$,
        and give their units.
        \vfill
  \item Use (b) to estimate the change in kinetic energy when a $50$~kg
        passenger gets in with the speed held fixed, and when the speed rises
        by $1$~m/s with the mass held fixed. Which change adds more energy?
        \vfill
\end{parts}
\end{prob}

\newpage
%==============================================================================
\section{Short proofs}
%==============================================================================
None of these needs a clever idea. Write down what you are assuming, justify
each step by naming the law or fact that licenses it, and end with a sentence
saying what you have shown.

\medskip
\noindent The laws of the dot product, from \S11.8 of Vector Unit, Part~2
(\emph{Orientation, the Dot Product, and Projection}), hold for all
$\vc{u},\vc{v},\vc{w}\in\R^{2}$ and all scalars $\alpha$:
\begin{center}
\renewcommand{\arraystretch}{1.4}
\begin{tabular}{@{}l l l@{}}
\textbf{D1} & $\ip{\vc{u}}{\vc{v}}=\ip{\vc{v}}{\vc{u}}$ & symmetry \\
\textbf{D2} & $\ip{\vc{u}}{\vc{v}+\vc{w}}=\ip{\vc{u}}{\vc{v}}+\ip{\vc{u}}{\vc{w}}$
            & additivity \\
\textbf{D3} & $\ip{\vc{u}}{\alpha\vc{v}}=\alpha\,\ip{\vc{u}}{\vc{v}}$
            & homogeneity \\
\textbf{D4} & $\ip{\vc{u}}{\vc{u}}=\norm{\vc{u}}^{2}\ge0$, with equality only
              for $\vc{u}=\vc{0}$ & positive definiteness \\
\end{tabular}
\end{center}

%------------------------------------------------------------------------------
\begin{prob}[A law of vector addition, in components]
Let $\vc{u}=[\,u_1,\,u_2\,]$, $\vc{v}=[\,v_1,\,v_2\,]$ and
$\vc{w}=[\,w_1,\,w_2\,]$ be vectors in $\R^{2}$. Working in components, prove
the law \textbf{A2} of \S8.4 of Vector Unit, Part~1:
\[
  (\vc{u}+\vc{v})+\vc{w}=\vc{u}+(\vc{v}+\vc{w}).
\]
Name the property of the real numbers that the proof rests on.
\end{prob}
\vfill

\newpage
%------------------------------------------------------------------------------
\begin{prob}[The real part, and a dot product in disguise]
Let $z=a+bi$ and $w=c+di$ with $a,b,c,d\in\R$.
\begin{parts}
  \item Prove that $z+\conj{z}=2\Rez(z)$ and $z-\conj{z}=2i\,\Imz(z)$.
        \vfill
  \item Prove that $z\conj{w}+\conj{z}w=2(ac+bd)$. So $z\conj{w}+\conj{z}w$ is
        always a real number, and it is twice the dot product of the vectors
        $[\,a,\,b\,]$ and $[\,c,\,d\,]$.
        \vfill
  \item Check (b) for $z=1+2i$ and $w=3-i$ by computing $z\conj{w}$ and
        $\conj{z}w$ separately and adding them.
        \vfill
\end{parts}
\end{prob}

\newpage
%------------------------------------------------------------------------------
\begin{prob}[The orthogonal piece, and an exact accounting]
Let $\vc{u},\vc{v}\in\R^{2}$ with $\vc{v}\neq\vc{0}$, and put
$c=\ip{\vc{v}}{\vc{u}}/\norm{\vc{v}}^{2}$, so that
$\proj{\vc{v}}{\vc{u}}=c\,\vc{v}$.
\begin{parts}
  \item Using only \textbf{D1}--\textbf{D4} --- no components --- prove that
        $\ip{\vc{v}}{\vc{u}-c\vc{v}}=0$.
        \hint{Write $\vc{u}-c\vc{v}=\vc{u}+(-c)\vc{v}$ and expand with
        \textbf{D2} and \textbf{D3}.}
        \vfill
  \item The \textbf{Pythagorean identity} (Vector Unit, Part~5, \S23) says that
        if $\ip{\vc{x}}{\vc{y}}=0$, then
        $\norm{\vc{x}+\vc{y}}^{2}=\norm{\vc{x}}^{2}+\norm{\vc{y}}^{2}$. Use it
        and (a) to prove that
        \[
          \norm{\vc{u}}^{2}
          =\norm{\proj{\vc{v}}{\vc{u}}}^{2}
           +\norm{\vc{u}-\proj{\vc{v}}{\vc{u}}}^{2}.
        \]
        \hint{Part (a) is about $\vc{v}$, but the identity needs
        $\ip{c\vc{v}}{\vc{u}-c\vc{v}}=0$. Move the $c$ out with \textbf{D1} and
        \textbf{D3}.}
        \vfill
\end{parts}
\end{prob}

\newpage
%------------------------------------------------------------------------------
\begin{prob}[A weighted norm]
For $\vc{v}=[\,v_1,\,v_2\,]\in\R^{2}$, define
\[
  \norm{\vc{v}}_{w}=\lvert v_1\rvert+2\lvert v_2\rvert .
\]
A function on $\R^{2}$ is a \textbf{norm} when it satisfies the three properties
of \S15.1 of Vector Unit, Part~3, for all $\vc{u},\vc{v}\in\R^{2}$ and all
scalars $\alpha$:
\begin{center}
\renewcommand{\arraystretch}{1.4}
\begin{tabular}{@{}l l l@{}}
\textbf{N1} & $\norm{\vc{v}}\ge0$, \ and \ $\norm{\vc{v}}=0$ if and only if
              $\vc{v}=\vc{0}$ & positive definiteness \\
\textbf{N2} & $\norm{\alpha\vc{v}}=\lvert\alpha\rvert\,\norm{\vc{v}}$
            & absolute homogeneity \\
\textbf{N3} & $\norm{\vc{u}+\vc{v}}\le\norm{\vc{u}}+\norm{\vc{v}}$
            & the triangle inequality \\
\end{tabular}
\end{center}
Prove that $\norm{\cdot}_{w}$ satisfies \textbf{N1}, \textbf{N2} and
\textbf{N3}. You may use two facts about real numbers $a$ and $b$:
$\lvert ab\rvert=\lvert a\rvert\,\lvert b\rvert$, and
$\lvert a+b\rvert\le\lvert a\rvert+\lvert b\rvert$.

\hint{For \textbf{N1}, a sum of two non-negative numbers is zero only when both
of them are.}
\end{prob}
\vfill

\newpage
%------------------------------------------------------------------------------
\begin{prob}[A span is a subspace]
Fix two vectors $\vc{u},\vc{v}\in\R^{2}$, and let
\[
  \spanop\{\vc{u},\vc{v}\}
  =\bigl\{\,\alpha\vc{u}+\beta\vc{v} : \alpha,\beta\in\R\,\bigr\}.
\]
\begin{parts}
  \item Prove that $\spanop\{\vc{u},\vc{v}\}$ is a subspace of $\R^{2}$: show
        that it is non-empty, closed under addition, and closed under scalar
        multiplication. You may use the fact that $0\vc{x}=\vc{0}$ for every
        vector $\vc{x}$, and these laws of \S8.4 of Vector Unit, Part~1, for all
        vectors $\vc{x},\vc{y},\vc{z}$ and scalars $\gamma,\delta$:
        \begin{center}
        \renewcommand{\arraystretch}{1.3}
        \begin{tabular}{@{}l l @{\qquad} l l@{}}
        \textbf{A1} & $\vc{x}+\vc{y}=\vc{y}+\vc{x}$ &
        \textbf{S1} & $\gamma(\vc{x}+\vc{y})=\gamma\vc{x}+\gamma\vc{y}$ \\
        \textbf{A2} & $(\vc{x}+\vc{y})+\vc{z}=\vc{x}+(\vc{y}+\vc{z})$ &
        \textbf{S2} & $(\gamma+\delta)\vc{x}=\gamma\vc{x}+\delta\vc{x}$ \\
        \textbf{A3} & $\vc{x}+\vc{0}=\vc{x}$ &
        \textbf{S3} & $\gamma(\delta\vc{x})=(\gamma\delta)\vc{x}$ \\
        \end{tabular}
        \end{center}
        \vfill
  \item Prove that if a vector $\vc{w}\in\R^{2}$ is orthogonal to both
        $\vc{u}$ and $\vc{v}$, then $\vc{w}$ is orthogonal to every vector in
        $\spanop\{\vc{u},\vc{v}\}$. Use the laws \textbf{D2} and
        \textbf{D3}.
        \vfill
\end{parts}
\end{prob}

\newpage
%==============================================================================
\section*{Where each question comes from}
%==============================================================================
Each question is modeled on the homework problems listed beside it. If one gives
you trouble, rework those problems first. HW1 is Homework~1 and HW2 is
Homework~2.

\begin{center}
\renewcommand{\arraystretch}{1.2}
\small
\begin{tabular}{@{}l l @{\qquad\qquad} l l@{}}
\toprule
\textbf{Question} & \textbf{Homework problems} &
\textbf{Question} & \textbf{Homework problems} \\
\midrule
A.1  & HW1 B.1, B.3       & C.1 & HW1 A.1--A.4, B.3, B.5 \\
A.2  & HW1 C.5, D.3       & C.2 & HW1 C.2--C.7 \\
A.3  & HW1 C.3            & C.3 & HW1 E.1, E.3; HW2 A.2, E.1 \\
A.4  & HW2 A.6            & C.4 & HW2 C.1--C.3 \\
A.5  & HW1 E.3            & C.5 & HW2 D.1, D.2, D.4 \\
A.6  & HW2 B.1, E.7       & C.6 & HW2 F.1, F.2, F.4, F.6 \\
A.7  & HW2 C.4            & D.1 & HW1 E.4; HW2 A.3 \\
A.8  & HW2 C.5            & D.2 & HW2 B.1, B.2 \\
A.9  & HW2 C.2, D.2       & D.3 & HW2 E.4 \\
A.10 & HW2 F.4, F.6       & D.4 & HW2 F.5 \\
\addlinespace
B.1  & HW2 A.1            & E.1 & HW1 B.2 \\
B.2  & HW1 C.2            & E.2 & HW1 C.4, D.1, D.2 \\
B.3  & HW1 C.5            & E.3 & HW1 E.7; HW2 D.4 \\
B.4  & HW1 E.3; HW2 A.2   & E.4 & HW2 B.3, E.2 \\
B.5  & HW1 E.4            & E.5 & HW2 C.4, C.5, D.5 \\
B.6  & HW2 C.1, C.2       &     & \\
B.7  & HW2 E.3            &     & \\
B.8  & HW2 F.4            &     & \\
\bottomrule
\end{tabular}
\end{center}

\vfill
\noindent\rule{\textwidth}{0.4pt}

\smallskip
\noindent{\small Questions about any problem here belong in the class Discord,
where everyone can see the answer. Questions that are specific to your own work
are better brought to office hours, Thursdays 3:00--4:00 PM in Kiely Hall 508.
Please email ahead.}

\end{document}
